\chapter{Why sharp-runtime Exists}
\index{sharp-runtime}\index{.NET compatibility}
\label{ch:sharp-runtime-overview}

This book's ecosystem chapter introduced \texttt{sharp-runtime} as the foundation layer every
XNA-facing CNA type quietly depends on. This chapter treats it as a project in its own right
--- one considerably larger, and considerably more broadly scoped, than its role inside CNA
alone might suggest.

\section{Scale}

At the audited revision, \texttt{sharp-runtime} contains 1{,}032 public headers plus five
private implementation headers, 217 implementation files, and 78{,}934 lines of production
C/C\texttt{++}. Its tests are larger
still: 102{,}342 lines, roughly 30\% more than production. The registered GoogleTest surface
contains 16{,}310 tests across 37 executables; the repository's continuation record reports
16{,}303 passing, one skipped, and six failing, rather than concealing those failures behind
an obsolete all-green count. This is a
considerably larger surface than ``just enough .NET to make XNA's API vocabulary work in
C\texttt{++}'' might suggest --- coverage genuinely extends to \texttt{System.Text.Json},
\texttt{System.Xml.Linq} / \texttt{XPath}, \texttt{System.Buffers.Text} / \texttt{Binary},
\texttt{System.Threading.Channels}, and \texttt{System.Net.WebSockets} / \texttt{Http} /
\texttt{Security}, well beyond a minimal XNA-support subset.

\section{Namespace structure}

The public headers live below two include roots. \texttt{System/} mirrors real .NET's
namespace tree directly:
\texttt{Collections} (with \texttt{Concurrent}, \texttt{Frozen}, \texttt{Generic},
\texttt{Immutable}, \texttt{ObjectModel}, \texttt{Specialized} sub-namespaces),
\texttt{IO} (with \texttt{Compression}, \texttt{Hashing}, \texttt{IsolatedStorage}),
\texttt{Net} (with \texttt{Http}, \texttt{Sockets}, \texttt{WebSockets}, \texttt{Security}),
\texttt{Text} (with \texttt{Json}, \texttt{RegularExpressions}, \texttt{Encodings}),
\texttt{Threading} (with \texttt{Channels}, \texttt{Tasks}), \texttt{Xml}, \texttt{Security},
\texttt{Globalization}, \texttt{Runtime}, and roughly 185 flat top-level types
(\texttt{Object}, \texttt{String}, \texttt{Exception}, \texttt{EventHandler},
\texttt{IDisposable}, \texttt{TimeSpan}, \texttt{DateTimeOffset}, \texttt{Delegate},
\texttt{Guid}, \texttt{Random}, \texttt{Uri}, \texttt{Version}, \texttt{Nullable},
\texttt{Span}, \texttt{Memory}, and more). A separate \texttt{SharpRuntime/}
namespace holds project-internal helpers that are not part of the .NET mirror at all ---
most importantly \texttt{SharpRuntimeHelper.hpp}, home to the type-alias table
Chapter~\ref{ch:design-philosophy} already introduced (\texttt{bytecs},
\texttt{intcs}, \texttt{Single}, and the rest).

There is no umbrella header. The on-disk include spelling follows the fully-qualified .NET
name mechanically: \texttt{System.IO.FileStream} becomes
\texttt{System/IO/FileStream.hpp}. CNA consequently consumes far more than the two headers
whose spelling begins with \texttt{SharpRuntime/}: its source directly includes 54 distinct
headers from this repository, 52 below \texttt{System/} and two below \texttt{SharpRuntime/}.
Those two are \texttt{SharpRuntimeHelper.hpp} and \texttt{Prop.hpp}; the latter exports
property-declaration macros used downstream even though the runtime's own modules do not use
it. Chapter~\ref{ch:sharp-components-consumption} maps these include paths to their owning
link components.

\section{Object and IDisposable: deliberately not two roots}

\cnaclass{System::Object} is not a universal root for this C\texttt{++} type system. It is an
abstract, vestigial compatibility type with a pure-virtual \texttt{GetTypeName()}, and only
\cnaclass{DateTime} and \cnaclass{DateTimeOffset} derive from it in the entire runtime.
\cnaclass{Stream} does not; exceptions derive from \texttt{std::exception}. Its
\texttt{GetHashCode()}, \texttt{Equals}, \texttt{ToString()}, and RTTI-backed
\texttt{GetType()} are useful to the few types that opt in, while the
\texttt{GetTypeNameHPP()} / \texttt{GetTypeNameCPP()} macros reduce their override
boilerplate. Reading those helpers as evidence that every ported class inherits
\cnaclass{Object} would import a .NET invariant the implementation intentionally does not
have.

\cnaclass{System::IDisposable} is a single pure interface method,
\texttt{virtual void Dispose() = 0}; its own documentation is explicit that C\texttt{++} has
no equivalent of C\#'s \texttt{using} statement, so callers either use RAII or call
\texttt{Dispose()} explicitly, and every implementation is expected to be idempotent.
It is a capability interface, not a second hierarchy root.

\subsection{String and Exception expose the same adaptation strategy}

\cnaclass{System::String} is a non-instantiable collection of static source-porting helpers;
both its constructor and destructor are deleted. Runtime strings remain \texttt{std::string},
also exported as \texttt{SharpRuntime::String}. Likewise,
\cnaclass{System::Exception} derives from \texttt{std::exception}, carries an
\texttt{exception\_ptr} for its inner exception and a default HRESULT, and returns an empty
stack trace because no managed stack is available. These are deliberate C\texttt{++}
representations of useful .NET behavior, not attempts to reproduce the CLR's object layout.

\section{EventHandler and MulticastAction: solving a problem C\# doesn't have}

\cnaclass{System::EventHandler<TEventArgs>} deliberately diverges from real .NET's own
design, and its own header comment states why directly: in C\#, \texttt{EventHandler<T>} is
only a delegate \emph{type} --- the actual multicast subscriber list and invocation mechanism
are supplied separately, by the compiler-generated code behind the \texttt{event} keyword.
C\texttt{++} has no \texttt{event} keyword, so \texttt{sharp-runtime}'s own
\cnaclass{EventHandler<T>} has to be both the type \emph{and} the subscriber-list
infrastructure at once: \texttt{operator+=}, a token-based \texttt{Add} / \texttt{Remove} pair,
\texttt{Clear()}, and \texttt{Raise(sender, e)}. \texttt{Raise} specifically invokes over a
\emph{snapshot} of the current handler list, not the live list --- a direct, deliberate guard
against a handler that removes itself (or another handler) from within its own invocation,
which previously surfaced as a real \texttt{std::bad\_function\_call} crash before this
snapshot discipline was adopted. A further mechanism, \texttt{SetReplayHook()}, exists
specifically to model a real XNA quirk this book has already touched on earlier ---
\cnaclass{NetworkSession.GamerJoined} (Chapter~\ref{ch:networking}) is expected to replay
already-happened state to a newly-added subscriber, not only notify about future events, and
this hook is the general mechanism that makes that pattern implementable once, rather than
reinvented per event.

\cnaclass{System::MulticastAction<Args...>} is a second, purpose-built multicast type built
for the same underlying reason: since \texttt{std::function}/lambda values have no C\#-style
target-plus-method equality, ordinary value comparison cannot implement C\#'s
\texttt{-=}-style single-subscriber removal, so this type uses the same
token-based-removal-plus-snapshot-invocation discipline as \cnaclass{EventHandler<T>} to solve
the identical problem for plain multicast callback fields that are not shaped like a
C\#-style public event. The motivating real-world case, named directly in this type's own
header comment, is re-parenting a scene-graph hierarchy: a node that resubscribes to a
different set of ancestor callbacks needs to drop its \emph{specific} old subscription without
clearing every other subscriber sharing the same field, which plain value comparison over
\texttt{std::function} objects cannot distinguish.

\begin{lstlisting}
// Worked example: token-based removal, the scenario this type exists for
// -- unsubscribing one specific handler without disturbing any other
// subscriber on the same field.
System::MulticastAction<float> onParentTransformChanged;
auto token = onParentTransformChanged.Add(
    [this](float deltaScale) { ApplyParentScale(deltaScale); });

// ... later, when this node is re-parented onto a different ancestor ...
onParentTransformChanged.Remove(token);          // only this subscription is dropped
token = newParent->onTransformChanged.Add(       // resubscribe to the new ancestor
    [this](float deltaScale) { ApplyParentScale(deltaScale); });
\end{lstlisting}

\section{Value types worth knowing before reading later chapters}

\cnaclass{System::TimeSpan}, the type underlying \cnaclass{GameTime} and
\cnaclass{SensorBase::TimeBetweenUpdates} throughout this book, stores its value internally
as 100-nanosecond ticks and carries one small, diagnostic-only surprise: static copy and move
counters exposed for test instrumentation. Both counters are now \texttt{std::atomic<int>};
loads, resets and increments use relaxed atomics because only each final count matters, not an
ordering relationship with the copied value. A downstream sanitizer suppression that still
describes their increments as a known \cnaclass{TimeSpan} race therefore documents an older
implementation, not the pinned runtime.

\cnaclass{Collections::Generic::List<T>} and \cnaclass{Dictionary<TKey,TValue>} both
implement .NET's fail-fast collection-modified-during-enumeration behavior via an internal
version counter checked on every enumeration step --- with one small, explicitly documented
deviation on \cnaclass{Dictionary}: its \texttt{Remove()} bumps the version counter, which
real .NET's own \texttt{Dictionary.Remove()} does not do, because \texttt{std::unordered\_map}'s
own iterator-invalidation guarantees on erase differ from the array-backed entry table real
.NET's dictionary uses internally --- a deliberate, reasoned adaptation to the underlying
C\texttt{++} standard library type, not an oversight. This is not just a two-type pattern:
every mutable generic collection this library ships --- \cnaclass{HashSet}, \cnaclass{LinkedList},
\cnaclass{Queue}, \cnaclass{Stack}, \cnaclass{SortedDictionary}, \cnaclass{SortedList},
\cnaclass{SortedSet}, and \cnaclass{OrderedDictionary} alongside the two above --- carries the
identical version-counter invariant, made project-wide as a deliberate late addition rather
than left as a two-type special case. \cnaclass{OrderedDictionary} specifically needed a
follow-up fix of its own even after that rollout: its \texttt{EnsureCapacity} method resized
the backing storage without bumping the version counter, meaning an enumerator already in
flight across a capacity-triggered resize would not detect the mutation the way every other
structural change on the same type correctly does --- closed as its own, narrower fix once
found, on the same underlying invariant.

\section{Concurrency primitives}

\cnaclass{Threading::Tasks::Task} implements a real, if scoped, async-task model:
\texttt{Wait()} rethrows a faulted task's exception (matching .NET's own
\texttt{Wait()} / \texttt{ThrowIfExceptional} contract), and \texttt{ContinueWith} uses
weak-pointer-based cycle avoidance specifically to avoid leaking a continuation chain rather
than spawning an extra OS thread per continuation. A detail worth stating plainly, since it
governs how a caller should reason about ordering: this port has no thread pool or scheduler
at all, so a continuation always runs synchronously and inline, either on whichever thread
completes the antecedent task or, if the antecedent is already complete by the time
\texttt{ContinueWith} is called, immediately on the calling thread --- effectively as if
\texttt{TaskContinuationOptions::ExecuteSynchronously} were always set, regardless of whether a
caller actually passes it. Most of the accepted \texttt{TaskContinuationOptions} filter bits
(\texttt{NotOnFaulted} / \texttt{NotOnCanceled} / \texttt{NotOnRanToCompletion} and their
\texttt{OnlyOnX} compositions) are genuinely honored; the rest
(\texttt{PreferFairness} / \texttt{LongRunning} / \texttt{AttachedToParent} /
\texttt{DenyChildAttach} / \texttt{HideScheduler} / \texttt{LazyCancellation}) are accepted only
for API-surface parity, since there is no scheduler or parent-task tracking for them to affect
in the first place.

\begin{lstlisting}
// Worked example: ContinueWith's continuation does NOT rethrow the
// antecedent's exception automatically the way Wait() does -- it must
// inspect the antecedent explicitly, and it runs synchronously, inline,
// with no extra thread spawned for it.
System::Threading::Tasks::Task loadTask = System::Threading::Tasks::Task::Run(
    [] { LoadLevelFromDisk(); });

System::Threading::Tasks::Task continuation = loadTask.ContinueWith(
    [](System::Threading::Tasks::Task antecedent)
    {
        if (antecedent.getIsFaultedProperty())
        {
            LogLevelLoadFailure();
            return;  // the antecedent's exception is NOT rethrown here automatically
        }
        FinishLevelSetup();
    });
\end{lstlisting}

\subsection{WhenAll cancellation and the external-future bridge}

\cnaclass{Task::WhenAll} and \cnaclass{Task::WhenAny} round out the async surface, and
\texttt{WhenAny} specifically was rebuilt to match real .NET's own zero-extra-thread design
(\texttt{TaskFactory.CommonCWAnyLogic}'s completion-registration strategy) rather than the
naive one-watcher-thread-per-input-task approach it started as. The current implementation
registers an inline callback on each input's shared state; the first callback wins an atomic
compare-exchange and later callbacks become cheap no-ops. No losing watcher is joined or
detached because the current path creates no watcher threads at all. \texttt{WhenAll} waits on
every input task to completion
(never short-circuiting on the first fault, matching real .NET) and, if one or more faulted,
rethrows the first one encountered in input order when the returned task is waited on --- a
deliberate simplification versus real .NET's \texttt{AggregateException} wrapping, consistent
with the same choice \cnaclass{Task::Wait} already makes for a single faulted task. If no input
faults but one is canceled, the action implementing \texttt{WhenAll} throws
\cnaclass{TaskCanceledException}; because that returned task has no associated cancellation
token, its constructor records the escape as \emph{Faulted}, not Canceled. Waiting still throws
the expected exception type, but the status flags remain a documented deviation from .NET.

\texttt{WhenAll} must distinguish that genuine canceled-input state from a task that merely
faulted with a directly thrown \cnaclass{TaskCanceledException}; it now checks the input task's
status instead of sniffing the exception type. The same category of defect previously existed
in the external-future bridge used by \texttt{TaskCompletionSource}: both
\texttt{TrySetCanceled()} and a caller's \texttt{SetException(TaskCanceledException)} settle a
promise with that exception type. A separate producer-set atomic flag now identifies only the
genuine cancellation path, preserving the caller's exception and Faulted state otherwise.
\texttt{WhenAny} needs neither rule: its outer \texttt{TaskT<Task>} always completes successfully
with the first terminal input, whose own fault/cancellation status the caller inspects.

\cnaclass{Threading::Thread} is a thin wrapper over \texttt{std::thread} that deliberately does
not start at construction --- a second call to \texttt{Start()} throws
\texttt{ThreadStateException}, matching real .NET's own one-shot-start contract exactly,
rather than silently allowing a thread object to be restarted the way a naive wrapper might.

\section{Permanent scope boundaries: what will never be ported, and why}

\texttt{sharp-runtime}'s own \texttt{CLAUDE.md} states its scope discipline directly, and it is
worth quoting rather than paraphrasing, since the wording distinguishes a real category this
book's own methodology depends on: ``Known permanent deviations (not bugs, not TODO).'' Five
areas are named explicitly, each with a stated reason rather than left as an unexplained gap ---
the delegate-tier deviation is the one already covered in full in
Chapter~\ref{ch:parity-philosophy} (the three-tier \texttt{Action} / \texttt{Delegate}/
\texttt{MulticastAction} split); the remaining four are covered here:

\begin{description}
  \item[Reflection] (\cnaclass{System::Type}, \cnaclass{System::Activator},
  \texttt{Enum.GetNames} / \texttt{GetValues}) --- ``completely out of scope. Stubs are the
  correct end state,'' not a gap awaiting engineering effort.
  \item[GC] (\texttt{System::GC}) --- mutating operations are callable no-ops and queries return
  documented sentinel values (for example zero, false, \texttt{NotApplicable}, or the constant
  three-generation maximum); memory is managed by RAII / \texttt{std::shared\_ptr} throughout,
  so there is no garbage collector for these calls to meaningfully forward to.
  \item[Serialization] (\texttt{[Serializable]}, \texttt{SerializationInfo}) --- ignored
  outright; not needed for game code.
  \item[P/Invoke and interop] --- out of scope entirely.
  \item[Cryptography] (symmetric/asymmetric ciphers, X.509 certificates, TLS ---
  \texttt{Aes*} / \texttt{RSA*} / \texttt{EC*} / \texttt{ChaCha20Poly1305} / \texttt{CryptoStream},
  \texttt{X509Certificates}, \texttt{SslStream}) --- out of scope by an explicit, dated decision
  (2026-07-07): correctly implementing this needs either a large new external dependency
  (OpenSSL/mbedTLS) or a hand-rolled, security-critical implementation, and neither is worth it
  for game code. This deviation is deliberately narrow, not a blanket ``no crypto'' rule:
  digest, authentication, and derivation APIs (\texttt{MD5} / \texttt{SHA*} / \texttt{HMAC} /
  \texttt{PBKDF2}) remain in scope and are already ported. HMAC keys and PBKDF2 passwords are
  sensitive material even though these APIs do not provide reversible confidentiality, X.509
  validation, or a TLS channel.
\end{description}

\subsection{Reflection's real stub, and the one contradiction it deliberately accepts}

\cnaclass{System::Type}'s own header comment is worth reading in full for what it says about
\emph{how} a permanent stub should behave, not just that it exists: constructed via
\texttt{Type::From<T>()}, backed by C++ RTTI's \texttt{std::type\_info} rather than any real
reflection metadata, its boolean predicates (\texttt{IsClass}, \texttt{IsValueType},
\texttt{IsAbstract}, \texttt{IsSealed}, \texttt{IsInterface}) each return a fixed value
regardless of what the actual type argument is --- because C++ RTTI cannot determine any of
them --- and the comment states plainly that these fixed values are \emph{not} mutually
consistent: \texttt{Type::From<int>()} reports both \texttt{IsClass()==true} and
\texttt{IsValueType()==false} simultaneously, a combination that would be self-contradictory for
a real .NET type (\texttt{int} is \texttt{IsValueType==true}, \texttt{IsClass==false}). This is
not an oversight left for a future fix --- the header says explicitly that these predicates
exist only so ported code calling \texttt{type.IsClass} still \emph{compiles}, not to answer the
question correctly, and that calling code must not branch on them to distinguish value types
from class types. \cnaclass{System::Activator} narrows the same permanent gap to exactly the
part that survives the loss of reflection: its real, callable
\texttt{CreateInstance<T>()} template constructs any default-constructible \texttt{T} at
compile time, while the reflection-based overloads real .NET exposes ---
\texttt{CreateInstance(Type)}, \texttt{CreateInstanceFrom} --- are absent entirely, not stubbed,
because \texttt{sharp-runtime} has no \texttt{System.Reflection} for them to resolve a
\texttt{Type} argument against in the first place.

\section{The build is part of the public architecture}

Since the 2026-08-10 modularization merge, \texttt{sharp-runtime} is not one static archive.
Its 41 module directories register 44 link components: 30 static libraries, 13 interface
components, and the \texttt{Xml.XPath} alias of \texttt{Xml}. Compatibility umbrellas named
\texttt{Core} and \texttt{Collections} coexist with the narrower components. The distinction
matters to CNA because an include path names a source-level dependency while a
\texttt{SharpRuntime::<Component>} target declares its link closure. The complete registry,
component selection rules, and CNA's old-checkout adapter are covered in
Chapter~\ref{ch:sharp-components-consumption}.

One platform-specific edge remains useful context here. On Android, the storage-path
implementation calls SDL3 path APIs, so a parent such as CNA must have created an
\texttt{SDL3::SDL3} target before adding the runtime. This conditional edge does not turn SDL
into a universal runtime dependency.

\section{Local and tracked CI gates have different scope}

The local gate configures and builds under strict shell failure handling and rejects compiler
warnings and errors. The modular test topology is finer-grained than the old
single-binary design: 36 component executables plus one integration executable feed GoogleTest
discovery. Each component test links its own component and declared test-only dependencies,
so an undeclared cross-module use can fail at link time instead of remaining hidden inside a
monolith. Python meta-tests and deliberately failing consumer fixtures add build-boundary
checks that ordinary positive unit tests cannot express.

The pinned repository also tracks \texttt{.github/workflows/components.yml}. On every push and
pull request it runs nine selective Ubuntu configurations, a full compatibility build, and a
separate Doxygen-warning job, with the build-job ceiling fixed at two. That proves a hosted
trigger and named Linux jobs exist; it does not prove Windows, macOS, Emscripten, or every
historical commit ran successfully. Older continuation prose saying no workflow existed was
superseded by the workflow now present in the pinned tree.

\subsection{Recorded failures and platform claims need scope}

The current continuation note records six failing and one skipped test; an older baseline that
claims every test passes is therefore historical, not authoritative. The same discipline
applies to platform reports: a past MinGW or Emscripten library build does not prove that all
37 current executables cross-compile, much less that they execute on their target. Native Linux
results establish neither a cross compiler's header and linker path nor runtime behavior on a
different operating system. Keeping those scopes separate is part of the evidence model, not
mere cautionary wording.
